\chapter{Formal Language of Behavioural States}
\label{ch:formal-language}

\section{Why This Chapter Must Avoid Fake Rigor}

The manuscript needs a formal language because later chapters discuss state
formation, state transition, and intervention design. But the language must be
honest about what kind of formalism it is. At present, the project does
\emph{not} have a fitted quantitative model with experimentally identified
coefficients for human daily behaviour. It therefore should not present
arbitrary formulas as if they were already established mathematics or settled
empirical law.

This chapter takes a stricter position. It provides:
\begin{enumerate}[nosep]
  \item a disciplined vocabulary for describing behavioural states;
  \item a heuristic scoring scaffold for comparing sources of resistance;
  \item qualitative monotonic assumptions about barriers and preparation;
  \item explicit evidence boundaries for what this notation does and does not
        claim.
\end{enumerate}

\begin{discussion}
The source support for this chapter comes from two directions. The
\textbf{textbook foundation} comes from dynamical-systems texts such as Eugene
Izhikevich's \emph{Dynamical Systems in Neuroscience}, Michael Brin and Garrett
Stuck's \emph{Introduction to Dynamical Systems}, and Steven Strogatz's
\emph{Nonlinear Dynamics and Chaos}. These sources justify using state-space,
stability, and transition language as mathematical scaffolding. The
\textbf{biological legitimacy} comes from later empirical chapters and their
sources, which show that structured state transitions and stability are not
mere metaphors. What they do \emph{not} supply is a calibrated scalar law for
everyday action failure.
\end{discussion}

\section{Four Levels of Formal Commitment}

Before introducing symbols, the chapter should state what sort of commitment
each symbol carries. Not every equation in a textbook has the same status.

\begin{enumerate}[nosep]
  \item \textbf{Descriptive notation.} Symbols that name analytically useful
        features, such as immediate reward pull or environmental reinforcement.
  \item \textbf{Heuristic scaffolds.} Simple summary expressions introduced for
        bookkeeping, ranking, or exposition.
  \item \textbf{Qualitative monotonic assumptions.} Claims such as ``more
        ambiguity usually raises transition difficulty'' that are plausible and
        structurally useful, but still not calibrated laws.
  \item \textbf{Unsupported if overstated.} Any claim of identified universal
        coefficients, uniquely correct functional form, or direct laboratory
        measurement for the chapter's everyday-action model.
\end{enumerate}

\begin{remark}
The practical rule for later chapters is simple: they may inherit the
notation, the scaffolds, and the monotonic assumptions, but they must not
quietly upgrade those elements into settled quantitative law.
\end{remark}

\section{State, Frame, and Descriptive Structure}

\begin{definition}[Behavioural state]
A \textbf{behavioural state} $s \in \mathcal{S}$ is a sustained mode of action
together with its associated posture, attention pattern, local environment, and
reward expectation.
\end{definition}

\begin{remark}
This is a modelling definition, not an empirical claim that every behavioural
state is a sharply separable biological object. Its purpose is to let the book
talk about persistence and transition without collapsing everything into the
single word ``motivation.''
\end{remark}

\begin{definition}[State profile]
For descriptive purposes, a behavioural state may be assigned a
\textbf{state profile}
\[
  \Pi(s) = \bigl(\Reward(s), \Switch(s), \Env(s)\bigr),
\]
where:
\begin{itemize}[nosep]
  \item $\Reward(s)$ denotes the immediate reward pull of remaining in $s$,
  \item $\Switch(s)$ denotes the cognitive cost of loading a competing task
        frame,
  \item $\Env(s)$ denotes environmental reinforcement of the current state.
\end{itemize}
\end{definition}

\begin{discussion}
The tuple $\Pi(s)$ is more defensible than prematurely forcing a universal
scalar. It says only that these three components are analytically useful and
should be named separately. That separation is already valuable for later case
analysis and intervention design, because it identifies different control
levers: reward structure, context loading, and environmental coupling.
\end{discussion}

\begin{remark}
At this stage, $\Pi(s)$ should be read as an \textbf{ordered descriptive
profile}, not as a measured coordinate system with fixed units. The chapter is
claiming that these dimensions are useful to separate, not that they have
already been standardized across persons, tasks, or laboratories.
\end{remark}

\section{A Heuristic Inertia Score, Not A Proven Law}

\begin{definition}[Heuristic inertia score]
If an analyst chooses a common normalization for the components of
$\Pi(s)$, a \textbf{heuristic inertia score} may be introduced:
\[
  \widehat{\FI}(s) = \alpha \Reward(s) + \beta \Switch(s) + \gamma \Env(s),
\]
with $\alpha,\beta,\gamma > 0$ chosen as modelling weights.
\end{definition}

\begin{remark}
The hat on $\widehat{\FI}$ matters. This expression is a bookkeeping device,
not an established law of behavioural dynamics. It is useful when one wants a
single comparison score, but the manuscript should not pretend that the weights
have been empirically identified or that the linear form is uniquely correct.
\end{remark}

\begin{discussion}
What is defensible here is the \textbf{structure of the scaffold}, not the
claim that the linear combination has already been validated as a universal
formula. Textbook and review sources justify using low-dimensional state
descriptions and stability language. They do not justify claiming that the
human cost of action obeys a known three-term linear law. The score is kept
because it is operationally helpful, but it is explicitly downgraded from
``formal law'' to ``heuristic model.''
\end{discussion}

\begin{remark}
The chapter could have avoided a scalar score entirely and remained purely
ordinal. The reason for retaining $\widehat{\FI}$ is pedagogical economy: later
chapters need a compact way to say that one state is more resistant than
another. The chapter therefore licenses the score as a \emph{comparison aid},
not as a deep empirical result.
\end{remark}

\section{Target States and Barriers}

\begin{definition}[Target state]
Given a current state $s_0$, a target state $s_1$ is a behavioural state whose
expected long-run value exceeds that of $s_0$ under the agent's stated goals
or declared evaluation rule.
\end{definition}

\begin{definition}[Activation barrier]
The \textbf{activation barrier} for transition $s_0 \to s_1$, written
$\DV(s_0 \to s_1)$, is a modelling construct representing how difficult it is
to leave the current regime and stabilize the target regime.
\end{definition}

\begin{remark}
The activation barrier is not assumed to be a literal physical energy, nor a
directly measured neural variable. It is an explanatory abstraction. The point
is to ask which factors make transition harder or easier, not to claim a known
laboratory unit for effort.
\end{remark}

\subsection{What We Can Defend}

The most defensible claim at this stage is qualitative rather than exact:
\begin{enumerate}[nosep]
  \item barriers tend to rise when current-state reinforcement is stronger;
  \item barriers tend to rise when the target state is ambiguous;
  \item barriers tend to fall when relevant preparation has already been done;
  \item barriers depend on current conditions \emph{and} transition history.
\end{enumerate}

\begin{definition}[Preparation and ambiguity]
Let $U(s_1)$ denote \textbf{target ambiguity}, meaning the unresolved
specification burden of entering $s_1$. Let $P$ denote \textbf{preparation},
meaning the extent to which entry-relevant work has already been performed
before the transition begins.
\end{definition}

\begin{remark}
\textbf{Monotone barrier scaffold.}
For modelling purposes, the transition difficulty may be treated as depending
on a function
\[
  \DV(s_0 \to s_1) \approx g\!\bigl(\widehat{\FI}(s_0), U(s_1), P, H\bigr),
\]
where $H$ summarizes transition-relevant history, and where $g$ is assumed to
be increasing in current-state resistance and ambiguity, and decreasing in
preparation.
\end{remark}

\begin{discussion}
This assumption is deliberately weaker and more defensible than the chapter's
earlier barrier-decomposition formula. It keeps the business logic that later
chapters need: if the present state is strongly reinforced, if the target state
is vague, and if preparation is weak, the transition becomes harder. But it no
longer pretends that the manuscript has identified a specific additive law with
known coefficients.
\end{discussion}

\begin{remark}
The symbol ``$\approx$'' should be read conservatively here. It does not mean
statistical fit or asymptotic approximation. It means only that the barrier
construct is being represented by a modelling dependence on those named inputs.
\end{remark}

\subsection{Ordinal Reading Before Cardinal Reading}

The safest use of the chapter's notation is ordinal:
\begin{itemize}[nosep]
  \item one state can be treated as more reinforcing than another,
  \item one target can be treated as more ambiguous than another,
  \item one preparation regime can be treated as better than another,
  \item one transition can be treated as harder than another.
\end{itemize}

This is already enough for most of the book's downstream claims. Intervention
design, case comparison, and decision-window reasoning usually need ranking and
direction of change more than exact measurement. A stronger cardinal reading
would require calibration work that this manuscript does not yet possess.

\subsection{What We Are Not Claiming}

This chapter does \emph{not} claim:
\begin{enumerate}[nosep]
  \item that $\widehat{\FI}$ has a universal physical scale;
  \item that the linear score is uniquely correct;
  \item that $\DV$ has already been measured for daily routines in a controlled
        quantitative programme;
  \item that the model distinguishes all psychologically relevant variables.
\end{enumerate}

These are real limitations. They should be stated now, because later chapters
depend on this notation. If the notation is honest about its status, later
extensions become more stable rather than less.

\section{Worked Example: Bed to Desk}

\begin{example}[A state-profile analysis rather than a fake calculation]
Consider the transition from scrolling in bed to seated reading at a desk.
\begin{enumerate}[nosep]
  \item $\Reward(s_0)$ is elevated because scrolling offers rapid novelty.
  \item $\Switch(s_0)$ is elevated because the reader must recover task
        context, locate materials, and define a first step.
  \item $\Env(s_0)$ is elevated because posture and location both stabilize
        rest-oriented behaviour.
  \item $U(s_1)$ is elevated if the study task is not pre-specified.
  \item $P$ rises when the desk is prepared, the text is opened in advance, and
        a first micro-task is written down.
\end{enumerate}
The model does not compute a true numerical barrier here. It gives a structured
explanation of why the transition feels expensive and why preparation can lower
that expense before the moment of action.
\end{example}

\begin{discussion}
This example is intentionally modest. It does not demonstrate predictive
success, parameter recovery, or empirical fit. Its purpose is narrower: to show
how the notation separates several otherwise collapsed complaints such as ``I
don't feel like it,'' ``I forgot where to start,'' and ``the room keeps pulling
me back into rest.''
\end{discussion}

\section{Why This Chapter Still Matters}

Removing fake rigor does not make the chapter weaker. It makes it usable.
Because the language is now more honest, later chapters can do three different
jobs without contradiction:
\begin{enumerate}[nosep]
  \item the empirical consciousness chapter can justify state and stability
        language biologically;
  \item the emergence chapter can explain how histories and repeated coupling
        shape later barriers;
  \item the advanced dynamics chapter can add richer mathematical structure
        without inheriting a false claim that the opening formula was already
        a proven law.
\end{enumerate}

\section{Evidence Boundaries}

\begin{itemize}[nosep]
  \item \textbf{Strong:} state, stability, and transition are legitimate
        scientific topics in biology and neuroscience.
  \item \textbf{Medium:} low-dimensional state-space language, attractor
        intuition, and monotone barrier reasoning are useful explanatory
        scaffolds for behavioural modelling.
  \item \textbf{Heuristic:} the scalar inertia score $\widehat{\FI}$ and the
        barrier function $g$ are modelling devices introduced for clarity and
        cross-chapter consistency.
  \item \textbf{Descriptive only:} the state profile $\Pi(s)$ and terms such as
        $\Reward$, $\Switch$, $\Env$, $U$, and $P$ should first be read as
        analytic categories before they are read as measurable variables.
  \item \textbf{Unsupported if overstated:} any claim that the exact linear
        coefficients or a universally valid additive barrier law have already
        been established for human everyday action.
\end{itemize}

\section{Recommended Reading}

\begin{readingbox}
\textbf{Foundation}
\begin{itemize}[nosep]
  \item Eugene Izhikevich, \emph{Dynamical Systems in Neuroscience}, for
        disciplined state-space and stability language.
  \item Michael Brin and Garrett Stuck, \emph{Introduction to Dynamical
        Systems}, for mathematical background on state evolution and local
        structure.
  \item Steven Strogatz, \emph{Nonlinear Dynamics and Chaos}, for threshold,
        stability, and regime-change intuition without overclaiming empirical
        fit.
\end{itemize}
\textbf{Modern Research and Review}
\begin{itemize}[nosep]
  \item Review literature on attractor and integrator networks in the brain
        (Nature Reviews Neuroscience, 2022, medium), for why state persistence
        and retained directional information are serious modelling ideas.
  \item \emph{Metastability demystified: the foundational past, the pragmatic
        present and the promising future} (Nature Reviews Neuroscience, 2024,
        medium), for why stable-but-revisable state descriptions are often more
        appropriate than crude rigid-versus-fluid binaries.
  \item \emph{Cortical connectivity, local dynamics and stability correlates of
        global conscious states} (Communications Biology, 2025, strong as a
        biological legitimacy source, not as direct validation of this
        chapter's scoring formula).
\end{itemize}
\textbf{Reading discipline}
\begin{itemize}[nosep]
  \item Use the textbooks to justify the modelling language.
  \item Use the review and biological sources to justify taking state,
        stability, and transition structure seriously.
  \item Do not read either group of sources as direct proof that this chapter's
        heuristic score has already been empirically calibrated.
  \item If later chapters introduce stronger formulas, they must explicitly say
        whether those formulas remain heuristic, merely monotone, or genuinely
        supported beyond what is established here.
\end{itemize}
\end{readingbox}

\begin{summarybox}
This chapter still gives the manuscript a formal language, but it now does so
honestly. What is kept is the useful structure: behavioural states, target
states, descriptive state profiles, heuristic resistance scoring, and monotone
barrier reasoning. What is removed is fake precision. The result is a more
stable foundation for the rest of the book: later chapters can build empirical,
emergent, and dynamical depth on top of a model that openly declares itself as
scaffold rather than pretending to be a finished quantitative law.
\end{summarybox}
